% from main.tex
%%%%%%%%%%  scan idx 497  |  folio 484  %%%%%%%%%%
Nous allons appliquer ces définitions aux $\lambda$-anneaux augmentés $K^{\bullet}(X)$ et $K^{\bullet}(Y)$, les anneaux d'augmentation $H^0(X,\mathbb Z)$ et $H^0(Y,\mathbb Z)$, étant binômiaux (V 2.8). Pour simplifier les notations, nous poserons, pour tout schéma $Z$, et tout $z\in K^{\bullet}(Z)$:
\[
\ch(z)=\ch(\widetilde c(z));
\qquad
\Todd(z)=\Todd(c(z)).
\]
Nous considérerons pour toute algèbre graduée $A$ le groupe $1+\widehat A^+$ comme contenu dans $\widehat A$. Enfin, si on a un homomorphisme $f:A\to B$, additif, et compatible aux graduations à un décalage près, on notera encore $f$ l'homomorphisme $\widehat A\to\widehat B$ qu'on en déduit.

\textbf{3.4.} Soit $X$ un schéma. Pour tout $\mathcal O_X$-Module localement libre de type fini $E$, on note $\check E$ la classe dans $K^{\bullet}(X)$ du dual $\check E$ de $E$. Il est clair que la fonction qui à $E$ associe $\check E$ est une fonction additive sur la catégorie des $\mathcal O_X$-Modules localement libres de type fini, et par suite définit un endomorphisme de $K^{\bullet}(X)$, qu'on notera $x\mapsto\check x$.

\textbf{Proposition 3.5.} Soit $X$ un schéma quasi-compact. Alors pour tout $x\in K^{\bullet}(X)$ et tout $i\geq0$, on a:
\[
c^i(\check x)=(-1)^i c^i(x).
\]
La relation à démontrer peut encore s'écrire :
\[
c_t(\check x)=c_{-t}(x),
\]
en posant : $c_t(x)=\sum_{i\geq0}c^i(x)t^i$.
